\chapter*{Linux 源码阅读索引}
\addcontentsline{toc}{chapter}{Linux 源码阅读索引}

本册不是 Linux 内核源码注释，但每个重要机制都应能落到真实源码。下面的路径以主线 Linux 内核常见目录为准；不同版本函数名会变化，阅读时以当前源码树为准。目标不是背路径，而是知道“这个概念在真实系统里由哪些结构和函数承载”。

\begin{longtable}{p{0.2\linewidth}p{0.34\linewidth}p{0.34\linewidth}}
\toprule
主题 & 关键路径 & 重点对象或函数 \\
\midrule
x86 启动 & \filepath{arch/x86/boot/}，\filepath{arch/x86/kernel/head\_64.S} & 实模式入口、长模式入口、早期页表 \\
内核初始化 & \filepath{init/main.c}，\filepath{arch/x86/kernel/setup.c} & \code{start\_kernel}、\code{setup\_arch}、initcall \\
系统调用入口 & \filepath{arch/x86/entry/entry\_64.S}，\filepath{arch/x86/entry/syscalls/syscall\_64.tbl} & \code{entry\_SYSCALL\_64}、\code{do\_syscall\_64} \\
x86 page fault & \filepath{arch/x86/mm/fault.c}，\filepath{arch/x86/mm/tlb.c} & CR2、error code、用户/内核 fault、TLB flush 和 shootdown \\
调度器 & \filepath{kernel/sched/core.c}，\filepath{kernel/sched/fair.c}，\filepath{kernel/sched/rt.c} & \code{\_\_schedule}、\code{enqueue\_task\_fair}、\code{pick\_next\_task} \\
进程创建退出 & \filepath{kernel/fork.c}，\filepath{kernel/exit.c} & \code{copy\_process}、\code{do\_exit}、\code{wait} \\
futex & \filepath{kernel/futex/} & \code{futex\_wait}、\code{futex\_wake}、PI futex \\
锁与 lockdep & \filepath{kernel/locking/}，\filepath{kernel/locking/lockdep.c} & qspinlock、mutex slow path、lock class graph \\
虚拟内存 & \filepath{mm/memory.c}，\filepath{mm/mmap.c}，\filepath{mm/rmap.c} & \code{handle\_mm\_fault}、\code{do\_mmap}、COW、VMA/PTE 生命周期 \\
页面回收 & \filepath{mm/vmscan.c}，\filepath{mm/swap.c}，\filepath{mm/oom\_kill.c} & \code{shrink\_lruvec}、kswapd、OOM killer \\
内核分配器 & \filepath{mm/page\_alloc.c}，\filepath{mm/slub.c}，\filepath{mm/vmalloc.c} & buddy、SLUB、\code{kmem\_cache\_alloc} \\
ELF 装载 & \filepath{fs/binfmt\_elf.c}，\filepath{fs/exec.c} & \code{load\_elf\_binary}、\code{create\_elf\_tables} \\
VFS 与 open & \filepath{fs/open.c}，\filepath{fs/namei.c}，\filepath{fs/read\_write.c} & \code{path\_openat}、\code{vfs\_read}、\code{vfs\_write} \\
page cache & \filepath{mm/filemap.c}，\filepath{mm/readahead.c}，\filepath{fs/fs-writeback.c} & \code{filemap\_fault}、readahead、writeback \\
ext4 与 jbd2 & \filepath{fs/ext4/}，\filepath{fs/jbd2/} & extent、mballoc、journal commit、recovery \\
块层 & \filepath{block/blk-mq.c}，\filepath{block/bio.c}，\filepath{block/mq-deadline.c} & \code{blk\_mq\_submit\_bio}、bio、request、tag \\
设备模型 & \filepath{drivers/base/}，\filepath{include/linux/device.h} & \code{device}、\code{driver}、\code{bus\_type}、kobject \\
PCI/NVMe/virtio & \filepath{drivers/pci/}，\filepath{drivers/nvme/host/}，\filepath{drivers/virtio/} & probe/remove、queue、interrupt、DMA \\
网络栈 & \filepath{net/core/}，\filepath{net/ipv4/}，\filepath{include/linux/skbuff.h} & \code{sk\_buff}、\code{tcp\_v4\_rcv}、\code{tcp\_sendmsg} \\
epoll & \filepath{fs/eventpoll.c} & \code{ep\_poll\_callback}、ready list、wait queue \\
namespace/cgroup & \filepath{kernel/nsproxy.c}，\filepath{kernel/cgroup/} & namespace clone、cgroup v2 controller \\
LSM/seccomp & \filepath{security/}，\filepath{kernel/seccomp.c} & \code{security\_*} hooks、BPF filter、seccomp actions \\
tracing/perf/eBPF & \filepath{kernel/trace/}，\filepath{kernel/events/}，\filepath{kernel/bpf/} & tracepoint、perf event、BPF program/map \\
panic/kdump & \filepath{kernel/panic.c}，\filepath{kernel/kexec\_core.c} & \code{panic}、oops、kdump、crashkernel \\
\bottomrule
\end{longtable}

\begin{rubricbox}
源码阅读最低标准：先找核心结构体，再找对象创建路径、状态转换路径、失败回滚路径和观测点。只搜索一个函数名，通常看不到系统设计。
\end{rubricbox}
